\chapter{Réalisation du système}

\section*{Introduction}
\addcontentsline{toc}{section}{Introduction}

Ce chapitre présente la réalisation pratique du système de supervision thermique, depuis l'infrastructure matérielle jusqu'à l'interface utilisateur, en passant par le pipeline complet de traitement. La première section justifie les choix technologiques. La deuxième présente le matériel utilisé. La troisième décrit la configuration de la plateforme embarquée et de la base de données temporelle, ainsi que les modèles pré-entraînés. La quatrième section détaille l'orchestration du pipeline de diagnostic. Les cinquième à neuvième sections présentent l'implémentation de chaque module de traitement : module d'acquisition Modbus et débitmètre, modèle Grey-Box, détection d'anomalies par Isolation Forest, classification hybride des causes, et prédiction de maintenance par régression Ridge. La dixième section décrit l'interface de visualisation, la onzième le système d'alertes, et la douzième les tests. Enfin, la treizième section présente le déploiement.

Les extraits de code complets des modules principaux sont fournis en annexes A1 à A10.
Un simulateur de données a été développé pour permettre la validation du pipeline ML en l'absence de données réelles, et est également utilisé pour l'évaluation des performances.

\section{Choix et justification du stack technique}

\subsection{Stack Backend}

Le Tableau~\ref{tab:stack_backend} présente les technologies du backend et leur rôle.

\begin{table}[htbp]
\centering
\caption{Stack Backend}
\label{tab:stack_backend}
\begin{tabular}{c p{0.22\textwidth} p{0.10\textwidth} p{0.40\textwidth}}
\toprule
\textbf{Logo} & \textbf{Technologie} & \textbf{Version} & \textbf{Rôle dans le projet} \\
\midrule
\fbox{\parbox[c][0.6cm][c]{0.6cm}{\centering\tiny [RPi]}} & Raspberry Pi OS & Debian 12 & Système d'exploitation embarqué ARM~64 \\
\fbox{\parbox[c][0.6cm][c]{0.6cm}{\centering\tiny [Py]}} & Python & 3.11 & Langage d'implémentation du backend et du pipeline ML \\
\fbox{\parbox[c][0.6cm][c]{0.6cm}{\centering\tiny [FA]}} & FastAPI + Uvicorn & 0.115.11 & Framework API REST asynchrone avec WebSocket natif \\
\fbox{\parbox[c][0.6cm][c]{0.6cm}{\centering\tiny [In]}} & InfluxDB & 2.x & Base de données time-series pour les mesures capteurs \\
\fbox{\parbox[c][0.6cm][c]{0.6cm}{\centering\tiny [pM]}} & pymodbus & 3.8.6 & Communication Modbus RTU sur bus RS485 \\
\fbox{\parbox[c][0.6cm][c]{0.6cm}{\centering\tiny [gp]}} & gpiozero & 2.0 & Interruptions matérielles GPIO pour le débitmètre \\
\bottomrule
\end{tabular}
\end{table}

\textbf{Raspberry Pi OS (Debian~12)}~: choix naturel pour une cible ARM~64. Il offre un écosystème Linux complet, une compatibilité native avec Python et les bibliothèques ML, et un support à long terme (LTS) indispensable pour un système industriel fonctionnant en continu.

\textbf{Python~3.11}~: langage principal du projet. Sa maturité dans le domaine du ML (scikit-learn, numpy), sa bibliothèque standard riche (smtplib, logging, pickle, asyncio), et sa disponibilité sur ARM en font le choix le plus adapté pour un pipeline alliant acquisition bas niveau et traitement algorithmique.

\textbf{FastAPI + Uvicorn}~: framework asynchrone performant. Le support natif des WebSocket permet la diffusion temps réel des données vers le frontend sans polling. La validation automatique des schémas via Pydantic et la documentation OpenAPI générée facilitent le développement et la maintenance.

\textbf{InfluxDB~2.x}~: base de données spécialisée time-series. Son modèle de données (measurement, tags, fields) correspond exactement aux besoins de l'acquisition à 1~Hz. Les requêtes Flux permettent d'agréger les données pour l'entraînement ML et la rétention configurable évite la saturation du stockage sur RPi.

\textbf{pymodbus~3.8.6}~: bibliothèque de référence pour le protocole Modbus en Python. Son implémentation asynchrone (AsyncModbusSerialClient) s'intègre nativement dans la boucle asyncio du backend, et sa gestion des timeouts et des erreurs de communication est robuste.

\textbf{gpiozero~2.0}~: bibliothèque d'abstraction GPIO. Elle simplifie la configuration des interruptions matérielles et fournit une API propre pour le comptage d'impulsions du débitmètre YF-S201, avec un filtrage anti-rebond intégré.

\subsection{Stack Frontend}

Le Tableau~\ref{tab:stack_frontend} présente les technologies du frontend.

\begin{table}[htbp]
\centering
\caption{Stack Frontend}
\label{tab:stack_frontend}
\begin{tabular}{c p{0.22\textwidth} p{0.10\textwidth} p{0.40\textwidth}}
\toprule
\textbf{Logo} & \textbf{Technologie} & \textbf{Version} & \textbf{Rôle dans le projet} \\
\midrule
\fbox{\parbox[c][0.6cm][c]{0.6cm}{\centering\tiny [Rct]}} & React & 18 & Bibliothèque de composants d'interface utilisateur \\
\fbox{\parbox[c][0.6cm][c]{0.6cm}{\centering\tiny [Vi]}} & Vite & 5 & Bundler et serveur de développement rapide \\
\fbox{\parbox[c][0.6cm][c]{0.6cm}{\centering\tiny [TW]}} & TailwindCSS & 3 & Framework CSS utilitaire pour le design responsive \\
\fbox{\parbox[c][0.6cm][c]{0.6cm}{\centering\tiny [Rch]}} & Recharts & 2 & Librairie de graphiques React pour les courbes temps réel \\
\fbox{\parbox[c][0.6cm][c]{0.6cm}{\centering\tiny [WS]}} & WebSocket & natif & Communication bidirectionnelle temps réel avec le backend \\
\bottomrule
\end{tabular}
\end{table}

\textbf{React~18}~: bibliothèque frontend mature. Son modèle de composants réactifs permet de mettre à jour individuellement chaque carte de température sans re-rendre toute la page, ce qui est essentiel pour un affichage à 1~Hz.

\textbf{Vite}~: bundler nouvelle génération. Son rechargement à chaud instantané accélère le développement, et sa configuration de build optimisée produit un bundle léger pour le déploiement sur le RPi.

\textbf{TailwindCSS}~: framework utilitaire. Il permet de construire rapidement une interface sombre et professionnelle sans écrire de CSS personnalisé, et son système de thème s'adapte parfaitement à l'affichage kiosk industriel.

\textbf{Recharts}~: librairie de graphiques pour React. Basée sur D3, elle fournit des composants prêts à l'emploi (LineChart, AreaChart) pour les courbes d'évolution de température et d'épaisseur de calcaire.

\textbf{WebSocket (natif)}~: le protocole WebSocket, intégré nativement aux navigateurs modernes, permet une communication full-duplex à faible latence avec le backend FastAPI, garantissant la mise à jour en temps réel des trois onglets.

\subsection{Outils ML}

Le Tableau~\ref{tab:stack_ml} présente les outils de Machine Learning.

\begin{table}[htbp]
\centering
\caption{Outils Machine Learning}
\label{tab:stack_ml}
\begin{tabular}{c p{0.22\textwidth} p{0.10\textwidth} p{0.40\textwidth}}
\toprule
\textbf{Logo} & \textbf{Technologie} & \textbf{Version} & \textbf{Rôle dans le projet} \\
\midrule
\fbox{\parbox[c][0.6cm][c]{0.6cm}{\centering\tiny [sk]}} & scikit-learn & 1.4.0 & IsolationForest, RandomForest, Ridge, metrics \\
\fbox{\parbox[c][0.6cm][c]{0.6cm}{\centering\tiny [np]}} & NumPy & 1.26.4 & Calculs matriciels, extraction de features, statistiques \\
\fbox{\parbox[c][0.6cm][c]{0.6cm}{\centering\tiny [mp]}} & Matplotlib & 3.8+ & Visualisations (ROC, t-SNE, histogrammes, courbes Ridge) \\
\fbox{\parbox[c][0.6cm][c]{0.6cm}{\centering\tiny [pk]}} & pickle & stdlib & Persistance et chargement des modèles entraînés \\
\bottomrule
\end{tabular}
\end{table}

\textbf{scikit-learn~1.4.0}~: bibliothèque ML de référence. Elle fournit les trois algorithmes utilisés (IsolationForest, RandomForestClassifier, Ridge) avec une API unifiée (fit/predict), ainsi que les métriques d'évaluation (f1\_score, r2\_score, roc\_curve). Sa maturité et sa documentation exhaustive justifient son choix.

\textbf{NumPy~1.26.4}~: fondement du calcul numérique. Utilisé pour la construction des vecteurs de caractéristiques, les statistiques glissantes (moyenne, variance) et les opérations matricielles du modèle Grey-Box.

\textbf{Matplotlib~3.8+}~: visualisation scientifique. Génère les figures d'évaluation (courbes ROC de l'Isolation Forest, projection t-SNE du Random Forest, courbes de régression Ridge, histogrammes des scores) sauvegardées dans \texttt{models/plots/}.

\textbf{pickle (stdlib)}~: sérialisation native Python. Permet de sauvegarder et recharger les modèles entraînés (RF, IF, scaler, label encoder) au format \texttt{.pkl} sans dépendance externe.

\subsection{Librairies Python}

Le Tableau~\ref{tab:libs_python} liste l'ensemble des librairies Python utilisées dans le projet.

\begin{table}[htbp]
\centering
\caption{Librairies Python utilisées}
\label{tab:libs_python}
\begin{tabular}{p{0.22\textwidth} p{0.10\textwidth} p{0.55\textwidth}}
\toprule
\textbf{Librairie} & \textbf{Version} & \textbf{Rôle dans le projet} \\
\midrule
fastapi & 0.115.11 & Framework API REST avec support WebSocket \\
uvicorn & 0.27.0 & Serveur ASGI pour l'exécution de FastAPI \\
pydantic & — & Validation des schémas de données et configuration \\
pymodbus & 3.8.6 & Client Modbus RTU pour la lecture des 6 capteurs \\
influxdb-client & 1.45.0 & Client InfluxDB pour l'écriture et la lecture des mesures \\
scikit-learn & 1.4.0 & Algorithmes ML (IF, RF, Ridge) et métriques d'évaluation \\
numpy & 1.26.4 & Calculs numériques et extraction de features \\
requests & — & Appels HTTP à l'API Telegram pour les alertes \\
gpiozero & 2.0 & Interruptions GPIO pour le comptage d'impulsions \\
matplotlib & 3.8+ & Génération des figures d'évaluation des modèles \\
smtplib & stdlib & Envoi d'emails via Gmail SMTP (passerelle) \\
logging & stdlib & Traçabilité et journalisation des événements système \\
pickle & stdlib & Sérialisation des modèles ML entraînés \\
pytest & — & Framework de tests unitaires et d'intégration \\
\bottomrule
\end{tabular}
\end{table}

Chaque librairie a été choisie pour sa pertinence dans son domaine~: \textbf{fastapi} pour ses performances asynchrones et son WebSocket natif, \textbf{pymodbus} comme seule bibliothèque Modbus asynchrone mature pour Python, \textbf{influxdb-client} pour l'intégration native avec InfluxDB~2.x, \textbf{scikit-learn} comme standard de facto du ML en Python, \textbf{gpiozero} pour sa simplicité d'utilisation des GPIO sur RPi, \textbf{smtplib} (stdlib) pour éviter une dépendance externe d'envoi d'email, \textbf{requests} pour sa simplicité d'appel aux API REST, \textbf{pytest} pour sa compatibilité avec asyncio et ses fixtures de test, et \textbf{matplotlib} pour ses capacités de visualisation scientifique. La combinaison de ces librairies couvre l'ensemble des besoins du projet~: acquisition bas niveau (pymodbus, gpiozero), traitement ML (scikit-learn, numpy), stockage temporel (influxdb-client), API temps réel (fastapi, uvicorn), notification (requests, smtplib), test (pytest) et visualisation (matplotlib).

\section{Matériel utilisé}

Cette section présente les différents composants matériels du système de supervision.

\begin{figure}[htbp]
  \centering
  \fbox{\parbox[c][5cm][c]{0.85\textwidth}{\centering\itshape [Photo 4.2.1 : Raspberry Pi 4 Model B (8 Go RAM)]}}
  \caption{Raspberry Pi 4 -- carte embarquée principale}
  \label{fig:rpi4}
\end{figure}

\begin{figure}[htbp]
  \centering
  \fbox{\parbox[c][5cm][c]{0.85\textwidth}{\centering\itshape [Photo 4.2.2 : Bloc d'alimentation du système]}}
  \caption{Bloc d'alimentation}
  \label{fig:alimentation}
\end{figure}

\begin{figure}[htbp]
  \centering
  \fbox{\parbox[c][5cm][c]{0.85\textwidth}{\centering\itshape [Photo 4.2.3 : Module de transmission Modbus RTU RS485]}}
  \caption{Module de transmission Modbus RTU RS485}
  \label{fig:module_rs485}
\end{figure}

\begin{figure}[htbp]
  \centering
  \fbox{\parbox[c][5cm][c]{0.85\textwidth}{\centering\itshape [Photo 4.2.4 : Convertisseur USB/RS485]}}
  \caption{Convertisseur USB/RS485}
  \label{fig:usb_rs485}
\end{figure}

\begin{figure}[htbp]
  \centering
  \fbox{\parbox[c][5cm][c]{0.85\textwidth}{\centering\itshape [Photo 4.2.5 : Capteur de température PT100]}}
  \caption{Capteur PT100}
  \label{fig:pt100}
\end{figure}

\begin{figure}[htbp]
  \centering
  \fbox{\parbox[c][5cm][c]{0.85\textwidth}{\centering\itshape [Photo 4.2.6 : Débitmètre YF-S201]}}
  \caption{Débitmètre YF-S201}
  \label{fig:yfs201}
\end{figure}

\begin{figure}[htbp]
  \centering
  \fbox{\parbox[c][5cm][c]{0.85\textwidth}{\centering\itshape [Photo 4.2.7 : Module d'acquisition et PT100 installés sur le moule]}}
  \caption{Installation du module et du PT100 sur le moule}
  \label{fig:installation_moule}
\end{figure}

\begin{figure}[htbp]
  \centering
  \fbox{\parbox[c][5cm][c]{0.85\textwidth}{\centering\itshape [Photo 4.2.8 : Câble HDMI et écran industriel]}}
  \caption{Câble HDMI et écran industriel}
  \label{fig:hdmi_ecran}
\end{figure}

\begin{figure}[htbp]
  \centering
  \fbox{\parbox[c][5cm][c]{0.85\textwidth}{\centering\itshape [Photo 4.2.9 : Prototype complet en production]}}
  \caption{Prototype complet en production}
  \label{fig:prototype}
\end{figure}

\section{Configuration de la plateforme}

\subsection{Configuration du système embarqué}

\begin{figure}[htbp]
  \centering
  \fbox{\parbox[c][5cm][c]{0.85\textwidth}{\centering\itshape [Photo 4.1 : Raspberry Pi~4 avec convertisseur USB/RS485 et câblage capteurs Modbus]}}
  \caption{Plateforme matérielle déployée}
  \label{fig:rpi_setup}
\end{figure}

La configuration de la carte ARM est entièrement scriptée dans \texttt{setup\_rpi.sh}, qui automatise les onze phases suivantes~:

\begin{enumerate}
  \item Durcissement SSH et paramètres réseau
  \item Installation d'InfluxDB~2.x (fichier \texttt{.deb} officiel + configuration)
  \item Installation de Node.js~20 via NodeSource
  \item Installation des dépendances Python dans un environnement virtuel (FastAPI, pymodbus, scikit-learn, gpiozero, influxdb-client)
  \item Configuration de l'affichage kiosk (Chromium en mode plein écran, rotation HDMI)
  \item Configuration du watchdog matériel et des cron de maintenance
  \item Installation du service backend (fichier \texttt{.service} systemd)
  \item Installation du service frontend (serveur Vite en mode production)
  \item Redémarrage des services et vérification d'intégrité
\end{enumerate}

La configuration réseau est statique, l'accès SSH est restreint aux clés, et le pare-feu limite les ports ouverts (8001~pour l'API, 8086~pour InfluxDB).

\subsection{Configuration de la base temporelle et modèles pré-entraînés}

InfluxDB est configuré avec un bucket unique \texttt{sensors} et une politique de rétention des données sur 30~jours. Le schéma de mesures stockées pour chaque cycle d'acquisition est~:

\begin{itemize}
  \item \textbf{Measurement \texttt{temperature}} : tags \texttt{group\_id}, \texttt{mold\_id}, fields \texttt{value}, \texttt{status}, \texttt{deviation}, \texttt{timestamp}.
  \item \textbf{Measurement \texttt{flow}} : tags \texttt{group\_id}, fields \texttt{lpm}, \texttt{timestamp}.
  \item \textbf{Measurement \texttt{diagnostic}} : tags \texttt{group\_id}, fields \texttt{cause}, \texttt{confidence}, \texttt{urgency}, \texttt{timestamp}.
  \item \textbf{Measurement \texttt{prediction}} : tags \texttt{group\_id}, fields \texttt{days\_remaining}, \texttt{ci\_lower}, \texttt{ci\_upper}, \texttt{timestamp}.
\end{itemize}

Ces données sont interrogées par le module d'évaluation des modèles pour le ré-entraînement périodique.

\subsubsection{Modèles pré-entraînés}

Les modèles ML entraînés sont persistés dans le répertoire \texttt{models/} :

\begin{center}
\begin{minipage}{0.70\textwidth}
\begin{verbatim}
models/
├── isolation_forest.pkl       # Modèle Isolation Forest
├── random_forest.pkl           # Random Forest (classification N2)
├── ridge_regression.pkl        # Ridge Regression (prédiction TTF)
├── scaler.pkl                  # StandardScaler (normalisation)
├── label_encoder.pkl           # Encodage des causes
├── training_report.json        # Rapport d'évaluation (métriques)
├── plots/                      # Figures générées par plots_evaluation.py
│   ├── isolation_forest_roc.png
│   ├── random_forest_confusion.png
│   ├── random_forest_tsne.png
│   ├── ridge_curves.png
│   └── feature_importances.png
└── .gitkeep
\end{verbatim}
\end{minipage}
\end{center}

Ces fichiers sont chargés au démarrage du backend. Le ré-entraînement quotidien (Section~4.4) met à jour les \texttt{.pkl} et régénère les figures dans \texttt{plots/}.

\section{Orchestration du pipeline de diagnostic}

\subsection{Boucle principale d'acquisition}

L'orchestrateur, implémenté dans le backend, exécute une boucle infinie à la cadence de 1~Hz. Chaque cycle suit la séquence codée ci-dessous en pseudocode~:

\begin{center}
\begin{minipage}{0.90\textwidth}
\begin{verbatim}
ALGORITHME 4.1 : Pipeline principal d'acquisition et de diagnostic

Entrée : Aucune (boucle continue)
Sortie : Mise à jour des variables de diagnostic et stockage

TANT QUE vrai FAIRE
    // Phase 1 — Acquisition (lecture bus + GPIO)
    readings ← modbus.read_all_sensors()
    flow_lpm ← flow_sensor.read_lpm()
    
    // Phase 2 — Historiques glissants
    POUR CHAQUE lecture DANS readings FAIRE
        historique[clé].append(lecture.temperature)
    FIN POUR
    
    // Phase 3 — Grey-Box (soft sensing)
    POUR CHAQUE lecture DANS readings FAIRE
        résultat_gb ← grey_box.compute(lecture, flow_lpm)
        stocker(résultat_gb)
    FIN POUR
    
    // Phase 4 — Extraction features + détection IF
    vecteur_if ← isoler_caractéristiques(historiques)
    SI vecteur_if ≠ None ALORS
        résultat_if ← if_model.predict(vecteur_if)
        // résultat_if.anomaly ∈ {True, False}
    SINON
        résultat_if ← {anomalie: Faux}
    FIN SI
    
    // Phase 5 — Classification (règles + RF)
    SI résultat_if.anomalie ALORS
        vecteur_rf ← construire_vecteur_10_dim(historiques)
        // Niveau 1 : règles physiques d'abord
        cause_règle ← appliquer_règles_physiques(vecteur_rf)
        SI cause_règle ≠ None ALORS
            diagnostic ← cause_règle
        SINON
            // Niveau 2 : RF pour les cas ambigus
            résultat_rf ← rf_model.predict(vecteur_rf)
            SI résultat_rf.confiance ≥ SEUIL_CONFIANCE ALORS
                diagnostic ← résultat_rf.cause
            SINON
                diagnostic ← "CAUSE_INDETERMINEE"
            FIN SI
        FIN SI
        // Enrichir avec les données AMDEC
        diagnostic ← enrichir_amdec(diagnostic)
    SINON
        diagnostic ← {cause: "NORMAL", confiance: 1.0}
    FIN SI
    
    // Phase 6 — Prédiction Ridge (périodique)
    résultat_ridge ← ridge.predict()
    
    // Phase 7 — Stockage et diffusion
    influx.écrire_mesures(readings, diagnostic, résultat_ridge)
    webSocket.diffuser({lectures, diagnostic, prédictions})
    
    // Phase 8 — Alertes
    SI résultat_if.anomalie ALORS
        alerting.envoyer(diagnostic.sévérité, diagnostic)
    FIN SI
    
    attendre(1 seconde)
FIN TANT QUE
\end{verbatim}
\end{minipage}
\end{center}

\subsection{Ré-entraînement des modèles}

Le ré-entraînement est orchestré automatiquement selon les modalités suivantes~:

\begin{itemize}
  \item \textbf{Isolation Forest} : ré-entraîné quotidiennement à 2~h du matin sur les 7~derniers jours de données. L'entraînement est déclenché par une tâche planifiée (\texttt{cron}) qui exécute le script \texttt{train\_models.py} sans option particulière.
  \item \textbf{Random Forest} : ré-entraîné simultanément après l'IF. Le script \texttt{train\_models.py} interroge InfluxDB pour obtenir les données, extrait les caractéristiques, applique l'auto-labeling à partir des règles physiques, et entraîne le modèle. Le jeu d'étiquettes est produit par la méthode \texttt{auto\_label()} du classifieur de causes.
  \item \textbf{Ridge Regression} : ré-entraîné également quotidiennement. Nécessite un minimum de 7~jours de données avant la première prédiction exploitable.
\end{itemize}

En complément du ré-entraînement planifié, un module de surveillance de la santé des modèles (\texttt{model\_evaluator}) évalue toutes les 600~cycles (~10~minutes) les performances courantes. Si trois évaluations consécutives détectent une dégradation (taux d'anomalie IF~\(>15\%\), F1-weighted RF~\(<0,75\)), un ré-entraînement exceptionnel est déclenché.

\subsection{Évaluation des modèles}

L'évaluation des modèles est réalisée par le script \texttt{train\_models.py} avec l'option \texttt{--eval}, qui produit un rapport de classification complet stocké dans \texttt{models/training\_report.json}. Les métriques suivantes y sont consignées~:

\begin{itemize}
  \item Isolation Forest : FPR (faux positifs), FNR (faux négatifs), score d'anomalie moyen.
  \item Random Forest : F1-score pondéré (weighted), accuracy, matrice de confusion. Ce F1 ne reflète que la performance du RF sur les cas ambigus qui lui sont soumis (cf.~Section~4.8), pas la performance du diagnostic global.
  \item Ridge Regression : RMSE, \(R^2\), intervalle de confiance à 90\%.
\end{itemize}

Les figures de performance (t-SNE des features, courbes ROC, matrices de confusion, courbe d'étalonnage Ridge) sont générées par \texttt{plots\_evaluation.py} et sauvegardées dans \texttt{models/plots/}.

\begin{table}[htbp]
\centering
\caption{Métriques d'évaluation des modèles ML}
\label{tab:metriques}
\begin{tabular}{p{0.25\textwidth}p{0.18\textwidth}p{0.18\textwidth}p{0.29\textwidth}}
\toprule
\textbf{Modèle} & \textbf{Métrique} & \textbf{Valeur} & \textbf{Source} \\
\midrule
Isolation Forest & FPR & \textit{À remplir} & \texttt{train\_models.py --eval} \\
Isolation Forest & FNR & \textit{À remplir} & \texttt{train\_models.py --eval} \\
Random Forest & F1-weighted & \textit{À remplir} & \texttt{train\_models.py --eval} \\
Random Forest & Accuracy & \textit{À remplir} & \texttt{train\_models.py --eval} \\
Ridge & RMSE & \textit{À remplir} & \texttt{train\_models.py --eval} \\
Ridge & \(R^2\) & \textit{À remplir} & \texttt{train\_models.py --eval} \\
Ridge & Largeur IC~90\% & \textit{À remplir} & \texttt{train\_models.py --eval} \\
\bottomrule
\end{tabular}
\end{table}

\section{Module d'acquisition -- Modbus et débitmètre}

\begin{figure}[htbp]
  \centering
  \fbox{\parbox[c][5cm][c]{0.85\textwidth}{\centering\itshape [Photo 4.2 : Convertisseur USB/RS485 et capteurs de température Modbus]}}
  \caption{Matériel d'acquisition -- bus Modbus}
  \label{fig:modbus_hardware}
\end{figure}

\subsection{Acquisition Modbus RTU}

Le module \texttt{ModbusManager} (Annexe~A1) gère la communication avec les six capteurs de température sur le bus RS485. Chaque capteur est interrogé séquentiellement via une requête de lecture de registre de maintien (code fonction~0x03). L'accès au bus est protégé par un verrou asynchrone (\texttt{asyncio.Lock}) pour garantir l'exclusion mutuelle.

La configuration matérielle du bus est~: 9600~bauds, 8 bits, sans parité, 1 bit de stop, timeout de 1~seconde. La température est codée sur un registre 16~bits avec un facteur d'échelle de 0,1~\si{\celsius} par bit.

Le mappage des adresses esclaves vers les groupes et moules est défini par dictionnaire dans \texttt{config.py} (\texttt{SENSOR\_MAP}) :

\begin{center}
\begin{minipage}{0.70\textwidth}
\begin{verbatim}
SENSOR_MAP = {
    (1, 1): (1, 0),   # Heater 1 — moule 1
    (2, 1): (2, 0),   # Heater 2 — moule 1
    (2, 2): (3, 0),   # Heater 2 — moule 2
    (3, 1): (4, 0),   # Heater 3 — moule 1
    (3, 2): (5, 0),   # Heater 3 — moule 2
    (4, 1): (7, 0),   # Heater 4 — moule 1
}
\end{verbatim}
\end{minipage}
\end{center}

Un mécanisme de reconnexion automatique est implémenté~: si toutes les lectures échouent, le système tente une reconnexion au bus après un délai de garde de 30~secondes.

\subsection{Acquisition du débitmètre YF-S201}

Le module \texttt{FlowSensor} (Annexe~A7) utilise une interruption matérielle sur la broche GPIO~17 pour compter les impulsions du débitmètre YF-S201. La relation débit-fréquence est~:

\[
Q~(\text{L/min}) = \frac{F~(\text{Hz})}{7,5}
\]

Un filtre anti-pics rejette les valeurs supérieures à 30~L/min (invraisemblables physiquement) en les remplaçant par la valeur nominale par défaut de 16,5~L/min.

\subsubsection{Structure du backend}

L'ensemble des modules backend est organisé selon l'arborescence suivante~:

\begin{center}
\begin{minipage}{0.70\textwidth}
\begin{verbatim}
backend/
├── main.py                   # Point d'entrée FastAPI + boucle acquisition
├── config.py                 # Configuration (SLU, POSITION_MAP, SENSOR_MAP, etc.)
├── config.py.example         # Exemple de configuration (sans secrets)
├── modbus_manager.py         # Module Modbus RTU (Annexe A1)
├── grey_box.py               # Modèle Grey-Box (Annexe A2)
├── anomaly_detector.py       # Isolation Forest (Annexe A3)
├── cause_classifier.py       # Classifieur N1/N2 (Annexe A4)
├── ridge_predictor.py        # Ridge Regression (Annexe A5)
├── alerting.py               # Système d'alertes (Annexe A6)
├── flow_sensor.py            # Débitmètre YF-S201 (Annexe A7)
├── modbus_scanner.py         # Utilitaire de scan du bus RS485
├── test_*.py                 # Tests unitaires (12 fichiers)
├── __init__.py
└── utils/
    └── features.py           # Extraction des caractéristiques
\end{verbatim}
\end{minipage}
\end{center}

Chaque module est conçu de manière indépendante avec une interface claire (méthodes \texttt{read()}, \texttt{predict()}, \texttt{compute()}, etc.), ce qui facilite les tests unitaires et la maintenance.

\section{Modèle Grey-Box -- Implémentation}

Le modèle Grey-Box (Annexe~A2) implémente les équations de la Section~3.4. À chaque cycle, pour chaque moule, il calcule~:

\begin{enumerate}
  \item Le \(\Delta T\) mesuré par soustraction de la température du moule à la consigne (45~\si{\celsius}).
  \item Le \(\Delta T_{\text{normal}}\) déterminé à la calibration.
  \item Le \(\Delta T_{\text{calcaire}}\) par différence.
  \item La puissance thermique \(Q\) à partir du débit massique et de \(C_p\).
  \item La résistance thermique \(R = \Delta T_{\text{calcaire}} / Q\).
  \item L'épaisseur \(e = R \cdot \lambda \cdot A\).
\end{enumerate}

La calibration initiale est effectuée automatiquement au démarrage en interrogeant la première valeur enregistrée dans InfluxDB pour chaque moule. En l'absence d'historique, une valeur par défaut de \(\Delta T_{\text{normal}} = 1,5~\si{\celsius}\) est utilisée.

Cinq niveaux d'urgence sont émis (cf.~Tableau~3.4), accompagnés de la dégradation relative en pourcentage.

\section{Détection d'anomalies -- Isolation Forest}

\subsection{Implémentation}

Le détecteur d'anomalies (Annexe~A3) utilise l'implémentation \texttt{IsolationForest} de scikit-learn avec les paramètres suivants~:

\begin{itemize}
  \item \texttt{n\_estimators} = 100
  \item \texttt{contamination} = 0,05 (taux d'anomalies attendu dans les données)
  \item \texttt{random\_state} = 42 (reproductibilité)
\end{itemize}

L'extraction des huit caractéristiques est effectuée sur une fenêtre glissante de 30~secondes (configurable via \texttt{FEATURE\_WINDOW\_SECONDS}). En deçà de la période de remplissage de la fenêtre, le détecteur retourne \texttt{None}.

Les données sont normalisées avant soumission au modèle (moyenne nulle, variance unitaire) via \texttt{StandardScaler}.

\subsection{Condition d'activation}

L'Isolation Forest est exécuté à chaque cycle d'acquisition. Si le score d'anomalie dépasse le seuil de 0,5, la variable \texttt{anomaly\_detected} est positionnée à \texttt{True} et déclenche l'étape de classification.

\section{Classification des causes -- Approche hybride N1/N2}

\subsection{Rappel de l'architecture}

Le classifieur de causes (Annexe~A4) implémente l'architecture à deux niveaux décrite en Section~3.6~:

\begin{itemize}
  \item \textbf{Niveau~1 (N1)} : règles physiques déterministes, exécutées en premier. Si une cause est identifiée avec certitude, le diagnostic est immédiat et le RF n'est pas sollicité.
  \item \textbf{Niveau~2 (N2)} : Random Forest, invoqué uniquement lorsque les règles physiques ne permettent pas de conclure (retour \texttt{None}).
\end{itemize}

\begin{figure}[htbp]
  \centering
  \fbox{\parbox[c][4cm][c]{0.85\textwidth}{\centering\itshape [Figure 4.1 : Diagramme de décision du classifieur -- règles (N1) puis RF (N2)]}}
  \caption{Diagramme de décision du classifieur N1/N2}
\end{figure}

\subsection{Règles physiques (N1)}

Les règles physiques examinent trois indicateurs principaux~:

\begin{itemize}
  \item \texttt{affected\_ratio} : proportion de moules du groupe présentant une température sous le seuil d'alerte.
  \item \texttt{sudden\_drop} : chute brutale de température \(>1~\si{\celsius}\) en 2~minutes.
  \item \texttt{flow\_status} : débit actuel par rapport au débit nominal (16,5~L/min).
\end{itemize}

L'arbre de décision des règles est le suivant~:

\begin{center}
\begin{minipage}{0.85\textwidth}
\begin{verbatim}
FONCTION physical_rules(affected_ratio, sudden_drop, flow_lpm, variance_T):
    SI affected_ratio > 0.8 ET sudden_drop = Vrai ET flow_lpm < 8.0 ALORS
        RETOURNER "HEATER_POMPE_HS"
    SINON SI affected_ratio > 0.8 ET sudden_drop = Vrai ALORS
        RETOURNER "HEATER_RESISTANCE_HS"
    SINON SI variance_T > 0.15 ET flow_lpm ≥ 14.0 ALORS
        RETOURNER "BULLES_AIR"
    SINON SI flow_lpm < 8.0 ET affected_ratio > 0.5 ALORS
        RETOURNER "NIVEAU_BAS_VANNE_PANNE"
    SINON
        RETOURNER None  // Cas ambigu → N2
    FIN SI
FIN FONCTION
\end{verbatim}
\end{minipage}
\end{center}

\subsection{Random Forest (N2)}

Le Random Forest reçoit un vecteur de dix caractéristiques (les huit de l'IF augmentées de \(R^2\) Grey-Box et de l'autocorrélation à retard~1). La configuration du modèle est~:

\begin{itemize}
  \item \texttt{n\_estimators} = 100
  \item \texttt{max\_depth} = 10
  \item \texttt{min\_samples\_leaf} = 5
  \item \texttt{class\_weight} = \texttt{'balanced'}
\end{itemize}

Le seuil de confiance minimal est fixé à 0,5. En dessous, la cause retournée est \texttt{CAUSE\_INDETERMINEE}.

\subsection{Remarque importante sur la performance du RF}

Il est essentiel de comprendre que le F1-score du Random Forest ne reflète pas la performance globale du système de diagnostic. En effet~:

\begin{itemize}
  \item Le RF ne traite qu'environ 30\% des anomalies détectées, correspondant aux cas que les règles physiques (N1) n'ont pas su classer -- c'est-à-dire les cas intrinsèquement les plus ambigus.
  \item Les 70\% de cas restants, résolus immédiatement par les règles, ne lui sont jamais soumis.
  \item Le déséquilibre des classes AMDEC fait que le RF voit principalement des cas rares ou difficiles, ce qui pénalise mécaniquement son F1 par rapport à un RF qui classifierait toutes les anomalies.
  \item Par conséquent, le F1-weighted du RF doit être interprété comme une mesure de sa performance sur les \emph{cas résiduels ambigus}, et non comme une évaluation du diagnostic complet.
\end{itemize}

\subsection{Étiquetage automatique pour l'entraînement}

Pour l'entraînement supervisé du RF, les étiquettes sont générées automatiquement par la méthode \texttt{auto\_label()} du classifieur. Cette méthode applique les règles physiques à des fenêtres de données historiques pour produire des paires (vecteur de caractéristiques, cause). Les fenêtres normales (aucune anomalie) sont étiquetées \texttt{NORMAL}. Cette approche auto-supervisée évite le coûteux étiquetage manuel tout en produisant un ensemble d'entraînement représentatif.

\section{Prédiction de maintenance -- Ridge Regression}

\subsection{Implémentation}

Le module \texttt{RidgePredictor} (Annexe~A5) implémente la régression Ridge avec validation croisée pour déterminer le paramètre de régularisation \(\alpha\). Les caractéristiques utilisées sont~:

\begin{itemize}
  \item Épaisseur de calcaire estimée (mm), lissée sur 7~jours.
  \item Taux de variation quotidien de l'épaisseur.
  \item Débit moyen sur la période.
  \item Température moyenne sur la période.
\end{itemize}

La prédiction produit une durée de vie résiduelle (Time-To-Failure, TTF) en jours.

\subsection{Intervalle de confiance par Bootstrap}

Pour chaque prédiction, mille échantillons bootstrap sont tirés avec remise dans l'historique d'entraînement. Chaque échantillon produit un modèle Ridge et une prédiction. L'intervalle de confiance à 90\% est formé par les 5\up{ème} et 95\up{ème} centiles de la distribution des mille prédictions.

Le ré-entraînement a lieu quotidiennement. La première prédiction exploitable est disponible après 7~jours de collecte de données.

\subsubsection{Structure des modèles et scripts ML}

Les modèles et scripts de Machine Learning sont organisés comme suit~:

\begin{center}
\begin{minipage}{0.70\textwidth}
\begin{verbatim}
ml/
├── train_models.py            # Entraînement et ré-entraînement
├── plots_evaluation.py        # Génération des courbes ROC, t-SNE, importances
├── simulate_data.py           # Simulateur pour validation ML
├── generate_rf_validation.py  # Jeu de validation spécifique RF
├── models/
│   ├── isolation_forest.pkl   # Modèle IF entraîné
│   ├── random_forest.pkl      # Modèle RF entraîné
│   ├── ridge_regression.pkl   # Modèle Ridge entraîné
│   ├── scaler.pkl             # Normalisation
│   ├── label_encoder.pkl      # Encodage des causes
│   ├── training_report.json   # Métriques d'évaluation
│   └── plots/                 # Figures générées
└── requirements.txt           # Dépendances ML
\end{verbatim}
\end{minipage}
\end{center}

Le simulateur (\texttt{simulate\_data.py}) peut être utilisé indépendamment pour générer des flux de données synthétiques, ce qui a permis de valider le pipeline ML avant le déploiement sur site.

\section{Interface utilisateur}

\subsection{Architecture frontend}

L'interface utilisateur (Annexe~A8 pour le composant WebSocket) est développée avec React~18 et Vite. Elle reçoit les données mises à jour à 1~Hz via une connexion WebSocket maintenue en permanence.

Les trois onglets fonctionnels sont~:

\begin{enumerate}
  \item \textbf{Supervision} : affiche les six moules répartis en quatre groupes selon la nomenclature du Tableau~\ref{tab:nomenclature}. Chaque carte indique la température, le statut (OK/alerte/critique), l'écart à la consigne, et la référence du moule.
  \item \textbf{Diagnostic intelligent} : synthèse du dernier diagnostic avec la cause, le score de confiance, le niveau d'urgence Grey-Box, et les actions AMDEC.
  \item \textbf{Maintenance prédictive} : graphiques d'évolution de l'épaisseur de calcaire, prédictions de durée de vie résiduelle avec intervalle de confiance.
\end{enumerate}

\subsubsection{Structure du frontend}

L'interface utilisateur est organisée selon l'arborescence suivante~:

\begin{center}
\begin{minipage}{0.70\textwidth}
\begin{verbatim}
frontend/
├── src/
│   ├── components/
│   │   ├── MoldCard.jsx        # Carte individuelle de moule
│   │   ├── Dashboard.jsx       # Onglet Supervision
│   │   ├── DiagnosticPanel.jsx # Onglet Diagnostic intelligent
│   │   ├── MaintenanceChart.jsx# Onglet Maintenance prédictive
│   │   ├── WebSocket.jsx       # Connexion WebSocket temps réel
│   │   ├── AlertBanner.jsx     # Bannière d'alerte
│   │   └── Layout.jsx          # Structure à onglets commune
│   ├── App.jsx                 # Point d'entrée React
│   ├── main.jsx                # Rendu ReactDOM
│   └── index.css               # Styles TailwindCSS
├── index.html
├── vite.config.js
├── tailwind.config.js
├── package.json
└── README.md
\end{verbatim}
\end{minipage}
\end{center}

Le composant \texttt{WebSocket.jsx} gère la connexion bidirectionnelle avec le backend FastAPI et distribue les données mises à jour aux trois onglets via le contexte React.

\begin{table}[htbp]
\centering
\caption{Nomenclature des moules par groupe}
\label{tab:nomenclature}
\begin{tabular}{p{0.18\textwidth}p{0.18\textwidth}p{0.28\textwidth}}
\toprule
\textbf{Heater} & \textbf{Moule} & \textbf{Référence} \\
\midrule
Heater~1 & Moule~1 & RBA\_178 \#03 \\
Heater~2 & Moule~1 & RBA\_178 \#01 \\
         & Moule~2 & RBA\_178 \#02 \\
Heater~3 & Moule~1 & RBA\_174 \#02 \\
         & Moule~2 & RBA\_174 \#06 \\
Heater~4 & Moule~1 & RBA\_174 \#08 \\
\bottomrule
\end{tabular}
\end{table}

\begin{figure}[htbp]
  \centering
  \fbox{\parbox[c][5cm][c]{0.85\textwidth}{\centering\itshape [Photo 4.3 : Capture d'écran de l'interface -- onglet Supervision]}}
  \caption{Capture d'écran -- onglet Supervision}
  \label{fig:ui_supervision}
\end{figure}

\begin{figure}[htbp]
  \centering
  \fbox{\parbox[c][5cm][c]{0.85\textwidth}{\centering\itshape [Photo 4.4 : Capture d'écran -- onglet Diagnostic intelligent]}}
  \caption{Capture d'écran -- onglet Diagnostic}
\end{figure}

\begin{figure}[htbp]
  \centering
  \fbox{\parbox[c][5cm][c]{0.85\textwidth}{\centering\itshape [Photo 4.5 : Capture d'écran -- onglet Maintenance prédictive]}}
  \caption{Capture d'écran -- onglet Maintenance}
\end{figure}

\section{Système d'alertes}

Le module d'alertes (Annexe~A6) assure l'envoi direct des notifications par deux canaux, sans intermédiaire~:

\begin{itemize}
  \item \textbf{Telegram} : envoi via l'API HTTP \texttt{api.telegram.org}. Un bot dédié publie dans un groupe opérateurs et un canal responsable. Le message inclut la cause, les moules affectés avec leur nomenclature, la criticité AMDEC, et les actions recommandées.
  \item \textbf{Email} : envoi via SMTP avec \texttt{smtplib} (Gmail). Destiné au responsable maintenance pour les alertes CRITICAL et SYSTEM.
\end{itemize}

Les trois niveaux de sévérité sont~:

\begin{itemize}
  \item \textbf{WARNING} (température 40--42~\si{\celsius}) : notification Telegram aux opérateurs.
  \item \textbf{CRITICAL} (température \(<40\)~\si{\celsius} ou anomalie confirmée) : Telegram + Email.
  \item \textbf{SYSTEM} (perte de communication bus) : Telegram + Email, urgence maximale.
\end{itemize}

\subsection{Pseudocode du module d'alertes}

Le module \texttt{alerting.py} implémente la logique d'envoi des notifications selon le pseudocode suivant~:

\begin{center}
\begin{minipage}{0.90\textwidth}
\begin{verbatim}
FONCTION envoyer_alerte(niveau, diagnostic)
    message ← formater_message(diagnostic)  // cause, moules, criticité, actions
    
    SI niveau ∈ {WARNING, CRITICAL, SYSTEM} ALORS
        // Canal Telegram (toujours)
        url ← "https://api.telegram.org/bot" + TOKEN + "/sendMessage"
        POST url AVEC {
            chat_id: TELEGRAM_CHAT_ID,
            text:    message,
            parse_mode: "HTML"
        }
    FIN SI
    
    SI niveau ∈ {CRITICAL, SYSTEM} ALORS
        // Canal Email (critique uniquement)
        serveur ← SMTP_SSL("smtp.gmail.com", 465)
        serveur.connexion(EMAIL_ADRESSE, EMAIL_PASSWORD)
        serveur.envoyer(
            de:    EMAIL_ADRESSE,
            à:     DESTINATAIRE_EMAIL,
            sujet: "[ALERTE " + niveau + "] " + diagnostic.cause,
            corps: message_texte
        )
        serveur.fermer()
    FIN SI
FIN FONCTION
\end{verbatim}
\end{minipage}
\end{center}

\subsection{Configuration SMTP et création du bot Telegram}

\paragraph{Configuration SMTP (Gmail)}~: Le fichier \texttt{config.py} contient les paramètres de connexion SMTP~:

\begin{center}
\begin{minipage}{0.70\textwidth}
\begin{verbatim}
EMAIL_SMTP_SERVER = "smtp.gmail.com"
EMAIL_SMTP_PORT = 465
EMAIL_ADDRESS = "supervision.thermique@gmail.com"
EMAIL_PASSWORD = "app_password_16_caracteres"
EMAIL_RECIPIENT = "responsable.maintenance@example.com"
\end{verbatim}
\end{minipage}
\end{center}

L'authentification utilise un mot de passe d'application Gmail (16~caractères) généré dans les paramètres de sécurité du compte Google. Le port 465 (SMTP over SSL) est utilisé.

\paragraph{Création du bot Telegram}~: Les étapes de création du bot sont~:
\begin{enumerate}
  \item Ouvrir Telegram et rechercher \texttt{@BotFather}.
  \item Envoyer la commande \texttt{/newbot} et choisir un nom d'affichage et un nom d'utilisateur (ex. \texttt{SupervisionThermiqueBot}).
  \item \texttt{BotFather} retourne un jeton API (ex. \texttt{123456:ABC-DEF1234ghIkl-zyx57W2v1u123ew11}) à copier dans \texttt{config.py}.
  \item Créer un groupe Telegram, ajouter le bot comme administrateur, puis envoyer un message dans le groupe.
  \item Récupérer l'ID du groupe via \texttt{https://api.telegram.org/bot<TOKEN>/getUpdates}.
  \item Copier l'ID du groupe (\texttt{chat\_id}) dans \texttt{config.py}.
\end{enumerate}

\begin{figure}[htbp]
  \centering
  \fbox{\parbox[c][4cm][c]{0.60\textwidth}{\centering\itshape [Photo 4.6 : Configuration du module d'alertes dans config.py]}}
  \caption{Configuration du module d'alertes}
  \label{fig:alert_config}
\end{figure}

\begin{figure}[htbp]
  \centering
  \fbox{\parbox[c][4cm][c]{0.60\textwidth}{\centering\itshape [Photo 4.7 : Email d'alerte reçu -- vue expéditeur]}}
  \caption{Email d'alerte reçu -- vue expéditeur}
  \label{fig:alert_email_sent}
\end{figure}

\begin{figure}[htbp]
  \centering
  \fbox{\parbox[c][4cm][c]{0.60\textwidth}{\centering\itshape [Photo 4.8 : Email d'alerte reçu -- vue destinataire]}}
  \caption{Email d'alerte reçu -- vue destinataire}
  \label{fig:alert_email_received}
\end{figure}

\section{Tests}

\subsection{Stratégie de test}

La stratégie de test suit le cycle en V défini au chapitre~2. Chaque module est testé unitairement, puis l'intégration des modules est vérifiée.

\subsection{Tests unitaires}

Trente-cinq tests unitaires couvrent les cinq modules principaux~:

\begin{itemize}
  \item \texttt{test\_modbus\_manager.py} : lecture, reconnexion, cas d'erreur.
  \item \texttt{test\_grey\_box.py} : calcul d'épaisseur, calibration, urgence.
  \item \texttt{test\_anomaly\_detector.py} : extraction features, prédiction IF.
  \item \texttt{test\_cause\_classifier.py} : règles physiques (N1), prédiction RF (N2), auto-label, cas ambigus (retour \texttt{None}).
  \item \texttt{test\_ridge\_predictor.py} : entraînement, prédiction, Bootstrap, intervalle de confiance.
\end{itemize}

\begin{table}[htbp]
\centering
\caption{Tests unitaires -- résultats}
\label{tab:tests_unitaires}
\begin{tabular}{p{0.25\textwidth}p{0.25\textwidth}p{0.15\textwidth}p{0.15\textwidth}}
\toprule
\textbf{Module} & \textbf{Test} & \textbf{Attendu} & \textbf{Obtenu} \\
\midrule
Modbus Manager & Lecture séquentielle 6 capteurs & 6 lectures & \textit{À remplir} \\
Modbus Manager & Reconnexion après timeout & Reconnexion dans 30~s & \textit{À remplir} \\
Grey-Box & Calcul épaisseur (delta=2°C, débit=16.5~L/min) & \(e \approx 0,24\)~mm & \textit{À remplir} \\
Grey-Box & Calibration jour~1 & delta normal = 1,5°C & \textit{À remplir} \\
Détecteur IF & Extraction features (30~s de données) & 8 features & \textit{À remplir} \\
Détecteur IF & Prédiction sur données normales & anomaly=False & \textit{À remplir} \\
Classifieur & Règle N1 -- pompe HS & cause = HEATER\_POMPE\_HS & \textit{À remplir} \\
Classifieur & Règle N1 -- bulles air & cause = BULLES\_AIR & \textit{À remplir} \\
Classifieur & Cas ambigu $\rightarrow$ N2 & RF sollicité & \textit{À remplir} \\
Classifieur & Confiance insuffisante & cause = INDETERMINEE & \textit{À remplir} \\
Classifieur & Auto-label sur fenêtre normale & étiquette = NORMAL & \textit{À remplir} \\
Ridge & Prédiction TTF & days\_remaining $>0$ & \textit{À remplir} \\
Ridge & Bootstrap IC~90\% & IC cohérent & \textit{À remplir} \\
\bottomrule
\end{tabular}
\end{table}

Un jeu de validation spécifique pour le Random Forest est généré par \texttt{tests/generate\_rf\_validation.py}. Ce script construit des vecteurs de caractéristiques 10D avec des causes connues, indépendamment des règles d'auto-labeling. Les cas ambigus (zones grises) sont spécialement conçus pour tester la vraie valeur ajoutée du RF. L'exécution avec \texttt{--eval} produit le rapport de classification sur ce jeu.

\subsection{Tests d'intégration}

Les tests d'intégration valident le fonctionnement de l'ensemble assemblé~:

\begin{table}[htbp]
\centering
\caption{Tests d'intégration -- résultats}
\label{tab:tests_integration}
\begin{tabular}{p{0.22\textwidth}p{0.30\textwidth}p{0.15\textwidth}p{0.15\textwidth}}
\toprule
\textbf{Scénario} & \textbf{Description} & \textbf{Attendu} & \textbf{Obtenu} \\
\midrule
Pipeline complet & Boucle 1~Hz : acquisition $\rightarrow$ stockage $\rightarrow$ diffusion & Cycle $<1$~s & \textit{À remplir} \\
Backend $\rightarrow$ InfluxDB & Mesures écrites et lues & Écriture + lecture OK & \textit{À remplir} \\
InfluxDB $\rightarrow$ RF retrain & Données historiques $\rightarrow$ modèle entraîné & Modèle produit & \textit{À remplir} \\
Backend $\rightarrow$ WebSocket & Diffusion temps réel vers frontend & Mise à jour 1~Hz & \textit{À remplir} \\
Alerte Telegram & Envoi notification anomalie & Message reçu & \textit{À remplir} \\
Alerte Email & Envoi notification critique & Email reçu & \textit{À remplir} \\
Reconnexion Modbus & Déconnexion physique du bus & Reprise auto 30~s & \textit{À remplir} \\
\bottomrule
\end{tabular}
\end{table}

\paragraph{Monkey patching des dépendances matérielles}~: Les tests d'intégration et unitaires utilisent la technique du \textit{monkey patching} pour simuler les dépendances matérielles absentes de l'environnement de développement. Avant chaque session de test, le fichier \texttt{tests/conftest.py} remplace les modules réels (\texttt{pymodbus}, \texttt{gpiozero}) par des doubles mockés via \texttt{unittest.mock.patch}. Par exemple, \texttt{ModbusManager} est patched pour retourner des lectures simulées sans nécessiter de bus RS485 physique, et \texttt{FlowSensor} est patched pour fournir une fréquence d'impulsions artificielle. Cette approche permet d'exécuter la suite complète de 35~tests unitaires et 7~tests d'intégration sur n'importe quelle machine de développement, sans matériel spécifique.

Les figures suivantes présentent les captures d'écran des sept scénarios d'intégration simulés~:

\begin{figure}[htbp]
  \centering
  \fbox{\parbox[c][4cm][c]{0.60\textwidth}{\centering\itshape [Photo 4.9 : Test intégration -- Pipeline complet (boucle 1~Hz)]}}
  \caption{Test d'intégration -- Pipeline complet}
  \label{fig:test_pipeline}
\end{figure}

\begin{figure}[htbp]
  \centering
  \fbox{\parbox[c][4cm][c]{0.60\textwidth}{\centering\itshape [Photo 4.10 : Test intégration -- Backend $\rightarrow$ InfluxDB]}}
  \caption{Test d'intégration -- InfluxDB}
  \label{fig:test_influxdb}
\end{figure}

\begin{figure}[htbp]
  \centering
  \fbox{\parbox[c][4cm][c]{0.60\textwidth}{\centering\itshape [Photo 4.11 : Test intégration -- InfluxDB $\rightarrow$ RF retrain]}}
  \caption{Test d'intégration -- Ré-entraînement RF}
  \label{fig:test_rf_retrain}
\end{figure}

\begin{figure}[htbp]
  \centering
  \fbox{\parbox[c][4cm][c]{0.60\textwidth}{\centering\itshape [Photo 4.12 : Test intégration -- WebSocket temps réel]}}
  \caption{Test d'intégration -- WebSocket}
  \label{fig:test_websocket}
\end{figure}

\begin{figure}[htbp]
  \centering
  \fbox{\parbox[c][4cm][c]{0.60\textwidth}{\centering\itshape [Photo 4.13 : Test intégration -- Alerte Telegram]}}
  \caption{Test d'intégration -- Alerte Telegram}
  \label{fig:test_telegram}
\end{figure}

\begin{figure}[htbp]
  \centering
  \fbox{\parbox[c][4cm][c]{0.60\textwidth}{\centering\itshape [Photo 4.14 : Test intégration -- Alerte Email]}}
  \caption{Test d'intégration -- Alerte Email}
  \label{fig:test_email}
\end{figure}

\begin{figure}[htbp]
  \centering
  \fbox{\parbox[c][4cm][c]{0.60\textwidth}{\centering\itshape [Photo 4.15 : Test intégration -- Reconnexion Modbus]}}
  \caption{Test d'intégration -- Reconnexion Modbus}
  \label{fig:test_modbus_reconnect}
\end{figure}

\subsection{Tests d'acceptation}

Les tests d'acceptation (prototype en laboratoire) valident la conformité au cahier des charges. Les résultats seront renseignés après la phase de tests en conditions réelles (Juin 2026).

\section{Déploiement}

Le déploiement cible un Raspberry Pi~4 (8~Go) exécutant Raspberry Pi~OS. La procédure est entièrement automatisée par \texttt{setup\_rpi.sh} (11 phases). Les services critiques sont gérés par systemd~:

\begin{itemize}
  \item \texttt{supervision-backend.service} : backend FastAPI + boucle d'acquisition.
  \item \texttt{supervision-frontend.service} : serveur Vite (production).
  \item \texttt{influxdb.service} : base de données temporelle.
\end{itemize}

Le système est conçu pour fonctionner en continu (24h/24, 7j/7) avec une sortie HDMI connectée à un écran industriel pour l'affichage kiosk de l'interface.

\section{Mise en œuvre et validation terrain}

\subsection{Procédure de mise en service}

La mise en œuvre du système sur site suit le processus suivant~:

\begin{enumerate}
  \item \textbf{Installation matérielle} : fixation des six capteurs PT100 sur les moules, câblage du bus RS485 entre les modules d'acquisition et le convertisseur USB/RS485, installation du débitmètre YF-S201 sur la canalisation d'eau, branchement de l'alimentation 5~V et de l'écran HDMI.
  \item \textbf{Déploiement logiciel} : exécution du script \texttt{setup\_rpi.sh} qui automatise l'ensemble des onze phases de configuration (cf.~Section~4.3).
  \item \textbf{Calibration} : lancement de la boucle d'acquisition en mode calibration. Le système enregistre les températures nominales pendant 30~minutes et calcule les \(\Delta T_{\text{normal}}\) pour chaque moule.
  \item \textbf{Validation fonctionnelle} : exécution du \texttt{demo\_cli.py} pour simuler chaque scénario de défaillance et vérifier la réaction du système.
\end{enumerate}

\subsection{Scénarios de validation par demo\_cli}

L'outil \texttt{demo\_cli.py} (Annexe~A9) a été utilisé pour valider le comportement du système face aux différents modes de défaillance. Il se connecte à l'API backend et change dynamiquement le mode de simulation. Les captures suivantes présentent chaque scénario~:

\begin{figure}[htbp]
  \centering
  \fbox{\parbox[c][4cm][c]{0.60\textwidth}{\centering\itshape [Photo 4.16 : Mode NORMAL -- températures stables 44--45\si{\celsius}]}}
  \caption{Scénario NORMAL -- fonctionnement nominal}
  \label{fig:demo_normal}
\end{figure}

\begin{figure}[htbp]
  \centering
  \fbox{\parbox[c][4cm][c]{0.60\textwidth}{\centering\itshape [Photo 4.17 : Mode GRADUAL\_DROP -- décroissance lente]}}
  \caption{Scénario GRADUAL\_DROP -- dérive progressive}
  \label{fig:demo_gradual}
\end{figure}

\begin{figure}[htbp]
  \centering
  \fbox{\parbox[c][4cm][c]{0.60\textwidth}{\centering\itshape [Photo 4.18 : Mode SUDDEN\_DROP -- chute brutale au 60\up{ème} cycle]}}
  \caption{Scénario SUDDEN\_DROP -- chute brutale}
  \label{fig:demo_sudden}
\end{figure}

\begin{figure}[htbp]
  \centering
  \fbox{\parbox[c][4cm][c]{0.60\textwidth}{\centering\itshape [Photo 4.19 : Mode HEATER\_FAIL -- tous les moules sous 42\si{\celsius}]}}
  \caption{Scénario HEATER\_FAIL -- panne de chauffage}
  \label{fig:demo_heater}
\end{figure}

\begin{figure}[htbp]
  \centering
  \fbox{\parbox[c][4cm][c]{0.60\textwidth}{\centering\itshape [Photo 4.20 : Mode PUMP\_FAIL -- chute globale + lectures erratiques]}}
  \caption{Scénario PUMP\_FAIL -- panne de pompe}
  \label{fig:demo_pump}
\end{figure}

\begin{figure}[htbp]
  \centering
  \fbox{\parbox[c][4cm][c]{0.60\textwidth}{\centering\itshape [Photo 4.21 : Mode NOISY -- 10\% de lectures erronées]}}
  \caption{Scénario NOISY -- bruit capteur}
  \label{fig:demo_noisy}
\end{figure}

\begin{figure}[htbp]
  \centering
  \fbox{\parbox[c][4cm][c]{0.60\textwidth}{\centering\itshape [Photo 4.22 : Auto-séquence -- enchaînement de tous les modes]}}
  \caption{Scénario AUTO -- séquence complète de démonstration}
  \label{fig:demo_auto}
\end{figure}

\section{Plan de maintenance}

Le plan de maintenance assure la pérennité et la fiabilité du système de supervision en production. Il couvre les aspects logiciels, matériels et données.

\subsection{Maintenance préventive}

\begin{table}[htbp]
\centering
\caption{Plan de maintenance préventive}
\label{tab:maintenance_prev}
\begin{tabular}{p{0.28\textwidth}p{0.15\textwidth}p{0.42\textwidth}}
\toprule
\textbf{Action} & \textbf{Périodicité} & \textbf{Description} \\
\midrule
Nettoyage écran industriel & Hebdomadaire & Dépoussiérage de l'écran et du boîtier RPi \\
Vérification connexions RS485 & Mensuelle & Contrôle visuel des borniers et serrage des vis \\
Sauvegarde base InfluxDB & Hebdomadaire & Export \texttt{influx backup} vers stockage externe \\
Mise à jour système & Trimestrielle & \texttt{apt update \&\& apt upgrade}, vérification compatibilité \\
Rotation logs & Mensuelle & Nettoyage des logs système (\texttt{journalctl}, \texttt{*.log}) \\
Test alertes Telegram/Email & Mensuelle & Envoi d'un message de test automatique \\
Nettoyage débitmètre YF-S201 & Trimestrielle & Démontage, détartrage et vérification du rotor \\
\bottomrule
\end{tabular}
\end{table}

\subsection{Maintenance corrective}

Les procédures suivantes sont documentées pour les interventions d'urgence~:

\begin{itemize}
  \item \textbf{Panne RPi} : remplacement par la carte de secours préconfigurée (image SD clone). Temps d'arrêt estimé~: 15~minutes.
  \item \textbf{Panne capteur PT100} : remplacement à l'identique, recalibration automatique au prochain cycle.
  \item \textbf{Panne convertisseur USB/RS485} : remplacement par unité de rechange, vérification des paramètres de communication (9600~bauds, 8N1).
  \item \textbf{Panne débitmètre} : remplacement YF-S201, vérification du facteur d'étalonnage (7,5~Hz par L/min).
  \item \textbf{Corruption base InfluxDB} : restauration depuis la dernière sauvegarde hebdomadaire.
\end{itemize}

\subsection{Surveillance de la santé du système}

Un module de watchdog logiciel intégré au backend vérifie en continu~:

\begin{itemize}
  \item La réactivité du bus Modbus (timeout $< 1$~s par capteur).
  \item L'intégrité de la connexion InfluxDB (ping toutes les 60~secondes).
  \item La fraîcheur des données frontend (WebSocket actif, mise à jour $< 2$~s).
  \item L'espace disque restant (alerte si $< 1$~Go).
  \item La température CPU du RPi (alerte si $> 70$~\si{\celsius}).
\end{itemize}

En complément, un script \texttt{healthcheck.sh} exécuté toutes les 5~minutes par cron vérifie ces métriques et notifie l'opérateur via le système d'alertes en cas de dégradation.

\section*{Conclusion}
\addcontentsline{toc}{section}{Conclusion}

Ce chapitre a présenté la réalisation complète du système. Le stack technique a été justifié, la plateforme configurée, et chaque module de traitement implémenté. L'architecture hybride N1/N2 du classifieur de causes permet de traiter efficacement la majorité des cas par des règles déterministes tout en réservant la puissance du Random Forest aux situations ambiguës. Les tests unitaires et d'intégration assurent la validation de chaque composant, conformément à la démarche du cycle en V. La mise en œuvre terrain via \texttt{demo\_cli} et le plan de maintenance garantissent le déploiement et l'exploitation en conditions réelles. Les extraits de code complets sont fournis en annexes A1 à A10.
